\section{Marks and kinds}

\begin{definition}[Mark]
  \label{def:mark}
  \lean{Obsmetrologia.Mark}
  \leanok
  A \emph{mark} is the arbitrary atomic primitive of the system: assembled from
  nothing, with no internal structure, it is where decomposition stops. Its
  on-disk surface form is a finite string, but the string is presentation, not
  identity: two marks may share a surface yet differ in role.
\end{definition}

\begin{proposition}[Kind]
  \label{thm:kind}
  \lean{Obsmetrologia.Kind}
  \leanok
  \uses{def:mark}
  A mark is a genus with two kinds, token and claim, dual and distinguished by
  position: in an utterance the token is first (identity-carrying, ``which
  one'') and the claim second (content-carrying, ``what is said'').
\end{proposition}

\begin{lemma}[Exactly two kinds]
  \label{lem:two-kinds}
  \lean{Obsmetrologia.Kind.two_kinds}
  \leanok
  \uses{thm:kind}
  Every kind is either \texttt{token} or \texttt{claim}. The cardinality two is
  forced, not chosen: one kind would collapse the four-cell codomain of
  articulating to two statuses, and three kinds would oversaturate it to eight
  cells, for which the system has no derivation.
\end{lemma}

\begin{lemma}[Token and claim are distinct]
  \label{lem:token-ne-claim}
  \lean{Obsmetrologia.Kind.token_ne_claim}
  \leanok
  \uses{thm:kind}
  The duality of the two kinds is real, not nominal: \texttt{token} $\neq$
  \texttt{claim}.
\end{lemma}
